\scanpage{43}
\mspage{17}

On considère que cette donnée équivaut : celle d'une
suite d'inclusions de drapeaux
\[
 d_1\subset d_2\mathinner{\cdotp\cdotp}\subset d_p\subset d_{p+1}=d'
 \qquad(p=\operatorname{long}.d')
\]
dont chacun est un segment initial strict dans
le \deleted{\unsure{précédent}}\inserted{suivant}.
Ceci donne donc \deleted{des \unsure{filt|fil}} une
suite de projections
\[
 \Sigma_{d'}=\Sigma_{d_{p-1}}\unsure{-|\longrightarrow}\Sigma_{d_p}
 \unsure{-|\longrightarrow}\mathinner{\cdotp\cdotp}\unsure{-|\longrightarrow}\Sigma_{d_2}
 \unsure{-|\longrightarrow}\Sigma_{d_1}
\]
dont on considère le mapping-cône simultané
qu'on désigne par
\[
 C_\Sigma(d_1,d_2,-,
 \underset{\substack{\Vert\\d'}}{d_{p-1}})
\]
ou aussi par \(C_{\Sigma,d,d'}\).
On veut déterminer les
morphismes de \unsure{transition|transitions} naturels entre
\[
 C_{\Sigma,d,d'}.
\]

\margin{\unsure{NB} \(d\) est l'ens des
éléments maximaux
des \(d_1,\mathinner{\cdotp\cdotp},d_p\)}

\margin{On reviendra sur cette
\uline{construction} \uline{plus tard}
Elle n'est pas tellement
utile finalement \(\cdots\)}

\uline{10.} Reprenons la description \unsure{ab ovo|Ab ovo}
(dans le cas \(\mathcal M=(\mathrm{Top})\), pour fixer les idées
--- c'était le contexte adopté au début de ces notes).

\deleted{On}
On a un espace topologique
% A wavy upper frame with two descending sides encloses the formula; its bottom is open.
\[
 \begin{array}{|c|}\hline
 (24)\quad\Sigma_*,
 \end{array}
\]
muni d'une relation d'ordre
\[
 \Sigma_{**}\subset\Sigma_*\times\Sigma_*
\]
ayant un graphe fermé\unsure{;|,} et tel que
\[
 \Sigma_{**}\xrightarrow{\operatorname{pr}_2=\mathrm{but}}\Sigma_*
\]
soit une \deleted{\unsure{plongement}}\inserted{immersion} --- plus

\margin{\textcircled{10}}
\margin{42}
